\documentclass[11pt,a4paper]{article}

\usepackage[T1]{fontenc}
\usepackage[utf8]{inputenc}
\usepackage{lmodern}
\usepackage{amsmath,amssymb,mathtools}
\usepackage{microtype}
\usepackage{xcolor}
\usepackage{hyperref}
\usepackage[margin=2.5cm]{geometry}
\usepackage{enumitem}
\usepackage{listings}

\hypersetup{
  colorlinks=true,
  linkcolor=blue!55!black,
  urlcolor=blue!55!black
}

\lstdefinestyle{fortran}{
  language=Fortran,
  basicstyle=\ttfamily\small,
  keywordstyle=\color{blue!60!black},
  commentstyle=\color{green!40!black},
  frame=single,
  columns=fullflexible,
  keepspaces=true,
  breaklines=true
}

\newcommand{\ket}[1]{\lvert #1\rangle}
\newcommand{\bra}[1]{\langle #1\rvert}
\newcommand{\avg}[1]{\left\langle #1\right\rangle}
\newcommand{\Hil}{\mathcal H}
\newcommand{\dd}{\mathrm d}

\title{Static Correlation Functions in \texttt{nssDMRG}\\
       \large Mathematical Conventions and Implementation Notes}
\author{}
\date{September 2026}

\begin{document}
\maketitle
\tableofcontents

\section{Purpose and scope}

The implementation in \texttt{DMRG\_MEASURE} explicitly constructs the
operator product $O_A(i)O_B(j)$ in the final renormalized basis and evaluates
its expectation value on the stored superblock state.  The principal
difficulties are that the two operators may belong to the same DMRG block or
to opposite blocks, and that fermionically odd operators require the correct
parity string.

The geometrical cases are
\begin{equation}
  i,j\in L, \qquad i,j\in R, \qquad i\in L,\;j\in R,
\end{equation}
plus the special local case $i=j$.

\section{Correlation functions and public interface}

The generic routine evaluates
\begin{equation}
  C_{AB}(i,j)
  =\bra{\Psi}O_A(i)O_B(j)\ket{\Psi},
\end{equation}
where $\ket{\Psi}$ is the stored superblock state in the target quantum-number
sector.  If \texttt{connected=.true.}, the result is instead
\begin{equation}
  C^{\mathrm c}_{AB}(i,j)
  =\avg{O_A(i)O_B(j)}
   -\avg{O_A(i)}\avg{O_B(j)}.
\end{equation}

The preferred key-based interface is
\begin{lstlisting}[style=fortran]
corr = Measure_Corr_DMRG(keyA,keyB,posA,posB)
corr = Measure_Corr_DMRG(keyA,keyB,posA,posB,connected=.true.)
\end{lstlisting}
In this form, the matrix, quantum-number shift, and operator type are obtained
together from \texttt{LIST\_OPERATORS}.

The lower-level interface is intended for conjugated or externally composed
operators:
\begin{lstlisting}[style=fortran]
corr = Measure_Corr_DMRG(OpA,dqA,OpB,dqB,posA,posB,typeA,typeB)
corr = Measure_Corr_DMRG(OpA,dqA,OpB,dqB,posA,posB,typeA,typeB,&
                         connected=.true.)
\end{lstlisting}
Here the caller is responsible for supplying consistent shifts and operator
types.

\section{Ordered products of an arbitrary number of operators}

The more general interface evaluates the ordered product
\begin{equation}
  C(i_1,\ldots,i_M)
  =\bra{\Psi}O_1(i_1)O_2(i_2)\cdots O_M(i_M)\ket{\Psi}.
\end{equation}
The input order is the algebraic order: positions need not be sorted and may
repeat.  For operators stored on the sites, the preferred call is
\begin{lstlisting}[style=fortran]
value = Measure_Product_DMRG(keys,positions)
\end{lstlisting}
where \texttt{keys(a)} identifies $O_a$ on site \texttt{positions(a)}.  The
matrix, shift, and fermionic grading are read from the corresponding site.
For factors constructed by the caller, use
\begin{lstlisting}[style=fortran]
value = Measure_Product_DMRG(Ops,dqs,types,positions)
\end{lstlisting}
with \texttt{dqs(:,a)} equal to $dq(O_a)$ and \texttt{types(a)} equal to
\texttt{"fermionic"} for an odd factor.  These metadata must describe the
supplied matrices; in particular, the shift of a conjugated operator changes
sign.

For a fixed-sector expectation value, the routine first applies the selection
rules
\begin{equation}
  \sum_{a=1}^{M}dq(O_a)=0,
  \qquad
  \sum_{a=1}^{M}p(O_a)=0\pmod 2.
\end{equation}
If either fails, the value is zero.  For each odd factor at position $i_a$,
the Jordan--Wigner representation contributes $P_k=(-1)^{N_k}$ on every
physical site $k<i_a$.  These parity factors and $O_a(i_a)$ are multiplied
in the original input order.  Consequently, all factors on a repeated site
are combined before that site enters the block-growth sequence, retaining
the distinction between, for example, $c_i^\dagger c_i$ and
$c_i c_i^\dagger$.

The resulting local matrices are transported separately through the left and
right blocks; an unused block carries the identity.  The final superblock
contraction includes transitions between quantum-number sectors when the
two block shifts are nonzero but compensate.  This construction therefore
covers products entirely within one block as well as products crossing the
cut, in both serial and MPI runs.

For example, the four-point correlator of a bond operator
$O_{4f}(i)=T^z_iT^z_{i+1}$ is obtained directly as
\begin{equation}
  \avg{O_{4f}(i)O_{4f}(j)}
  =\avg{T^z_iT^z_{i+1}T^z_jT^z_{j+1}},
\end{equation}
by passing the four keys and positions.  No intermediate operator
$O_{4f}$ needs to be stored.  In particular, overlapping bonds retain
the stated operator order.

\section{Quantum-number shifts}

For every operator the convention is
\begin{equation}
  dq(O)=Q_{\mathrm{out}}-Q_{\mathrm{in}},
\end{equation}
or, equivalently,
\begin{equation}
  O\ket{q}\longmapsto\ket{q+dq(O)}.
\end{equation}
For a product of Abelian operators,
\begin{equation}
  dq(O_AO_B)=dq_A+dq_B.
\end{equation}

Since the expectation value is evaluated within a fixed total sector,
\begin{equation}
  \avg{O_AO_B}\ne 0
  \quad\Longrightarrow\quad
  dq_A+dq_B=0.
\end{equation}
Failure of this condition is a selection rule rather than an invalid input;
the routine therefore returns zero.

For example,
\begin{align}
  dq(S^+)+dq(S^-)&=0,\\
  dq(c^\dagger_{\uparrow})+dq(c_{\uparrow})
    &=(+1,0)+(-1,0)=(0,0).
\end{align}
Conversely, in a particle-number eigenstate,
\begin{equation}
  \avg{c^\dagger_{i\uparrow}c^\dagger_{j\downarrow}}=0,
\end{equation}
because its total shift is $(+1,+1)$.  Such an operator may still represent a
transition between different sectors; the present routine computes an
expectation value within one fixed sector.

\section{Fermionic grading}

The quantum-number shift does not determine the fermionic sign structure.
One must also specify the $\mathbb Z_2$ grading
\begin{equation}
  p(O)=N_f(O)\bmod 2,
\end{equation}
where $N_f(O)$ is the number of elementary $c$ or $c^\dagger$ factors in $O$.
Examples are
\begin{align}
  p(c)=p(c^\dagger)&=1,\\
  p(cn)&=1,\\
  p(n)=p(c^\dagger c)&=0,\\
  p(S^+)=p(c^\dagger_\uparrow c_\downarrow)&=0,\\
  p(\Delta)=p(c_\downarrow c_\uparrow)&=0,\\
  p(P)=p((-1)^N)&=0.
\end{align}
Composition obeys
\begin{equation}
  p(O_AO_B)=p(O_A)+p(O_B)\pmod 2.
\end{equation}

In the current implementation, \texttt{type="fermionic"} denotes an odd
operator.  For expectation values in a state of definite fermion parity,
the total operator must be even.  Thus even--even and odd--odd products are
allowed; even--odd products vanish.  An odd--odd product is even overall, but
requires a parity string.

The shift and grading are independent metadata.  For example, the pair
operator $\Delta^\dagger=c^\dagger_\uparrow c^\dagger_\downarrow$ is even but
has a nonzero quantum-number shift.

\section{Superblock state and sector decomposition}

The stored state is expanded in the product basis of the final left and right
blocks:
\begin{equation}
  \ket{\Psi}
  =\sum_{\alpha,\beta}\Psi_{\alpha\beta}
   \ket{\alpha}_L\otimes\ket{\beta}_R.
\end{equation}
If $\ket{\alpha}_L\in\Hil_L(q)$ and the total quantum number is $Q$, the
compatible right state belongs to $\Hil_R(Q-q)$.  Consequently,
\begin{equation}
  \ket{\Psi}=\bigoplus_q\ket{\Psi_q},
  \qquad
  \ket{\Psi_q}\in\Hil_L(q)\otimes\Hil_R(Q-q).
\end{equation}

During measurement initialization, \texttt{LI} and \texttt{RI} store the
block states associated with each superblock sector.  The inverse maps
\texttt{Lmap} and \texttt{Rmap} convert a full block-state index into its
index inside a sector.  They are used to construct rectangular operator
blocks efficiently when $dq\ne0$.

\section{Local case: equal positions}

For $i=j$, the product is first formed in the local Hilbert space, preserving
the requested order:
\begin{equation}
  O_{AB}(i)=O_A(i)O_B(i).
\end{equation}
It is then embedded into the growing block and transported to the final
renormalized basis.  No Jordan--Wigner string is introduced in this case.
The ordering is essential; for example,
\begin{equation}
  c_i^\dagger c_i=n_i,
  \qquad
  c_ic_i^\dagger=1-n_i.
\end{equation}

\section{Transport through block growth}

Let $U_n$ map the enlarged block basis at step $n$ to the retained basis.
An operator transforms as
\begin{equation}
  \widetilde O^{(n)}=U_n^\dagger O^{(n)}U_n.
\end{equation}
After the next site is introduced,
\begin{equation}
  O^{(n+1)}
  =\widetilde O^{(n)}\otimes X_{n+1},
\end{equation}
where $X_{n+1}$ is either a requested local factor or the identity.  In
schematic form,
\begin{equation}
  O^{(n)}
  \longrightarrow U_n^\dagger O^{(n)}U_n
  \longrightarrow
  \left(U_n^\dagger O^{(n)}U_n\right)\otimes X_{n+1}.
\end{equation}
The actual tensor-product ordering is chosen consistently with left/right
block growth and with the OBC/PBC convention.

The generic constructor \texttt{Build\_Product\_Block\_DMRG} performs this
operation.  It maps physical positions to growth indices using
\texttt{b2gMap}, starts with the first nontrivial factor entering the block,
inserts identities at all unrelated growth steps, and continues until the
final block basis is reached.

\section{Why same-block products must be composed before truncation}

Suppose that both operators belong to the same block.  The correct
representation of their product in the retained basis is
\begin{equation}
  U^\dagger O_AO_BU.
\end{equation}
Projecting the operators separately and multiplying afterwards gives instead
\begin{equation}
  (U^\dagger O_AU)(U^\dagger O_BU)
  =U^\dagger O_A\underbrace{UU^\dagger}_{\mathcal P}O_BU.
\end{equation}
Since $U$ is generally rectangular,
\begin{equation}
  UU^\dagger=\mathcal P\ne I,
\end{equation}
and therefore
\begin{equation}
  U^\dagger O_AO_BU
  \ne
  (U^\dagger O_AU)(U^\dagger O_BU)
\end{equation}
in general.  The local factors must consequently be joined before all later
truncations.

\section{Parity-even operators in the same block}

For parity-even operators at distinct positions,
\texttt{Build\_Corr\_Block\_DMRG} delegates to the generic product
constructor.  The procedure is:
\begin{enumerate}[label=\arabic*.]
  \item map each physical position to its block-growth index;
  \item determine which nontrivial factor enters first;
  \item build that first factor;
  \item apply each subsequent $U_n$;
  \item insert identities at intermediate sites;
  \item insert the other factor when its site enters;
  \item continue to the final block basis.
\end{enumerate}

The resulting operator is either $O_{AB}^L$ or $O_{AB}^R$, and its average is
evaluated as
\begin{equation}
  \bra{\Psi}O_{AB}^L\otimes I_R\ket{\Psi}
  \quad\text{or}\quad
  \bra{\Psi}I_L\otimes O_{AB}^R\ket{\Psi}.
\end{equation}

\section{Odd fermionic operators in the same block}

For two odd operators at $i<j$, the tensor-product representation used by
\texttt{nssDMRG} is
\begin{equation}
  O_iO_j
  \longrightarrow
  (O_iP_i)P_{i+1}\cdots P_{j-1}O_j,
  \qquad P_k=(-1)^{N_k}.
\end{equation}
For example,
\begin{equation}
  c_i^\dagger c_j
  \longrightarrow
  (c_i^\dagger P_i)P_{i+1}\cdots P_{j-1}c_j.
\end{equation}
The factor list passed to the generic constructor is therefore
\begin{equation}
  \left[O_iP_i,\,P_{i+1},\,\ldots,\,P_{j-1},\,O_j\right].
\end{equation}

If the requested order has the first argument on the right, the string is
still assembled in spatial left-to-right order.  Exchanging the two odd
endpoints once produces
\begin{equation}
  O_A(i)O_B(j)=-O_B(j)O_A(i),
  \qquad i\ne j,
\end{equation}
and the final operator is multiplied by $-1$.

\section{Operators on opposite sides of the cut}

For $i\in L$ and $j\in R$, parity-even operators give
\begin{equation}
  C_{AB}(i,j)
  =\bra{\Psi}O_A^L(i)\otimes O_B^R(j)\ket{\Psi}.
\end{equation}
The two factors are independently transported to the final bases of their
blocks.  This expectation value cannot be evaluated by a same-block average,
because neither tensor factor is the identity.  It requires an explicit
superblock contraction.

\subsection{Sector transitions}

The left operator maps
\begin{equation}
  O_L:\Hil_L(q_{\mathrm{in}})\longrightarrow\Hil_L(q_{\mathrm{out}}),
  \qquad
  q_{\mathrm{out}}=q_{\mathrm{in}}+dq_L.
\end{equation}
Writing $q=q_{\mathrm{out}}$ gives
\begin{equation}
  q_{\mathrm{in}}=q-dq_L.
\end{equation}
Thus the rows of the filtered operator belong to the output sector $q$, and
its columns belong to the input sector $q-dq_L$.  The corresponding
rectangular block is
\begin{equation}
  A_q=
  O_L\big|_{\Hil_L(q-dq_L)\rightarrow\Hil_L(q)}.
\end{equation}

Conservation of the total quantum number requires
\begin{equation}
  dq_R=-dq_L,
\end{equation}
so that the right rectangular block performs the compatible transition.
For each output sector $q$, the contribution is
\begin{equation}
  C_q
  =\bra{\Psi_q}A_q\otimes B_q\ket{\Psi_{q-dq_L}},
\end{equation}
and the full correlation is
\begin{equation}
  C_{AB}(i,j)=\sum_q C_q.
\end{equation}

\subsection{Matrix form of the contraction}

For a fixed sector, reshape the superblock vector as a matrix
$\Psi_{\alpha\beta}$, where $\alpha$ and $\beta$ are left- and right-block
indices.  Then
\begin{equation}
  \Phi_{\alpha\beta}
  =\sum_{\alpha',\beta'}
   A_{\alpha\alpha'}B_{\beta\beta'}
   \Psi_{\alpha'\beta'},
\end{equation}
or, subject to the vectorization convention,
\begin{equation}
  \Phi=A\Psi B^{\mathsf T}.
\end{equation}
The expectation value is
\begin{equation}
  C=\sum_{\alpha,\beta}
  \Psi_{\alpha\beta}^*\Phi_{\alpha\beta}.
\end{equation}

The implementation applies the two factors successively and never constructs
the full Kronecker matrix $A\otimes B$.  This substantially reduces memory
usage.

\subsection{Distributed contraction}

In MPI mode the superblock vector is distributed among processes.  The same
mathematical operation is performed in two data layouts:
\begin{equation}
  \Psi
  \xrightarrow{B}X
  \xrightarrow{\text{MPI transpose}}X^{\mathsf T}
  \xrightarrow{A}\Phi^{\mathsf T}
  \xrightarrow{\text{MPI transpose}}\Phi.
\end{equation}
Each process evaluates its local scalar product
\begin{equation}
  C_r=\bra{\Psi_r}\Phi_r\rangle,
\end{equation}
and the total value is obtained through
\begin{equation}
  C=\sum_r C_r.
\end{equation}

\section{Odd fermionic string across the cut}

For $i\in L$, $j\in R$, and $i<j$,
\begin{equation}
  O_iO_j
  =(O_iP_i)P_{i+1}\cdots P_L
   P_{L+1}\cdots P_{j-1}O_j.
\end{equation}
The string is split into
\begin{align}
  O_L^{\mathrm{string}}
  &=(O_iP_i)P_{i+1}\cdots P_L,\\
  O_R^{\mathrm{string}}
  &=P_{L+1}\cdots P_{j-1}O_j.
\end{align}
The correlator becomes
\begin{equation}
  C_{AB}(i,j)
  =\bra{\Psi}
   O_L^{\mathrm{string}}\otimes O_R^{\mathrm{string}}
   \ket{\Psi}.
\end{equation}
For adjacent sites at the cut, this reduces to the convention already used
for hopping terms,
\begin{equation}
  (c_L^\dagger P_L)\otimes c_R.
\end{equation}
If the requested operator order is right-to-left, one exchange of odd
endpoints produces the additional minus sign.

\section{Connected correlations}

For \texttt{connected=.true.}, the disconnected contribution is subtracted
after the full two-point function has been computed.  A one-point average can
be nonzero only for an even operator with zero shift:
\begin{equation}
  dq(O)=0,
  \qquad
  p(O)=0.
\end{equation}
If either condition fails, the one-point function vanishes by symmetry, and
the implementation avoids evaluating it.

\section{Spin--spin wrapper}

The spin wrapper uses
\begin{equation}
  \mathbf S_i\cdot\mathbf S_j
  =S_i^zS_j^z
   +\frac12S_i^+S_j^-
   +\frac12S_i^-S_j^+.
\end{equation}
Accordingly,
\begin{equation}
  \avg{\mathbf S_i\cdot\mathbf S_j}
  =\avg{S_i^zS_j^z}
   +\frac12\avg{S_i^+S_j^-}
   +\frac12\avg{S_i^-S_j^+}.
\end{equation}
Each term separately satisfies $dq_A+dq_B=0$.  Spin operators are
fermionically even, including the electronic realization
$S^+=c^\dagger_\uparrow c_\downarrow$, and therefore require no parity
string.

\section{Spin-orbital-resolved density correlations}

The density wrapper returns the matrix
\begin{equation}
  C_{ab}(i,j)=\avg{n_{ia}n_{jb}},
  \qquad a,b=1,\ldots,N_{so},
\end{equation}
where
\begin{equation}
  N_{so}=N_{\mathrm{spin}}N_{\mathrm{orb}},
  \qquad
  a=i_{\mathrm{orb}}+(i_{\mathrm{spin}}-1)N_{\mathrm{orb}}.
\end{equation}
For each spin-orbital,
\begin{equation}
  n_a=c_a^\dagger c_a,
  \qquad dq(n_a)=0,
  \qquad p(n_a)=0.
\end{equation}
Consequently, no fermionic string is needed.  In connected mode,
\begin{equation}
  C^{\mathrm c}_{ab}(i,j)
  =\avg{n_{ia}n_{jb}}-\avg{n_{ia}}\avg{n_{jb}}.
\end{equation}
The total-density correlation can be recovered as
\begin{equation}
  \avg{N_iN_j}
  =\sum_{a,b=1}^{N_{so}}C_{ab}(i,j),
  \qquad
  N_i=\sum_a n_{ia}.
\end{equation}

\section{Complete measurement flow}

The generic measurement proceeds as follows:
\begin{enumerate}[label=\arabic*.]
  \item Initialize the measurement state and load the blocks, rotation
        matrices, sector decomposition, and superblock vector.
  \item Validate positions and the dimensions of the two shifts.
  \item Apply the selection rules
        \begin{equation*}
          dq_A+dq_B=0,
          \qquad
          p_A+p_B=0\pmod 2.
        \end{equation*}
  \item If $i=j$, compose the two operators locally in the requested order.
  \item If both positions belong to one block, build the complete product
        before subsequent truncations.  Insert the parity string for two odd
        endpoints.
  \item If the positions belong to opposite blocks, build the two final block
        operators separately, split the fermionic string at the cut when
        necessary, filter all rectangular sector transitions, and evaluate
        $\bra{\Psi}O_L\otimes O_R\ket{\Psi}$.
  \item If requested, subtract the disconnected one-point contribution.
\end{enumerate}

The two independent pieces of physical metadata are therefore
\begin{equation}
  \boxed{dq(O)}
  \qquad\text{and}\qquad
  \boxed{p(O)}.
\end{equation}

\end{document}
